% Analysis / Discussion -- drop-in fragment. Requires \usepackage{booktabs}.

\section{Analysis}
\label{sec:analysis}

\subsection{The Threshold Is an Averaging Effect}
\label{sec:analysis-statistical}
The flatness of \texttt{min\_time} and score across trace count
(Section~\ref{sec:results-threshold}) localises the mechanism. The per-trace
cache-collision signal---the small timing difference the attack exploits---is a
fixed quantity that does not depend on how many traces are collected. What
changes with count is the variance of each group mean $T_{i,j}[d]$: with more
traces per ciphertext-difference group, the correct final-round relation
separates from the noise floor more reliably. Recovery is therefore governed by
whether the group-mean estimates are stable enough to rank the correct relation
first, which is a smooth, probabilistic function of count---hence the
psychometric curve rather than a hard boundary.

\subsection{Why Go Needs About Five Times More Traces}
\label{sec:analysis-crosslang}
The same attack, on the same key, recovers with $\approx 5\times$ more traces in
Go than in C. Because the recovery logic is common, the difference originates in
collection: Go's measured timings are a noisier estimate of the same underlying
signal, so more traces are needed to average that noise down to the level where
the correct relation wins. The garbage-collector ablation
(Section~\ref{sec:results-gc}) sharpens this. If GC pauses were the dominant
noise source, disabling the collector would move the threshold substantially;
instead it moves $N_{50}$ by only $\approx 8\%$ and leaves the gap to C intact.
The residual noise is therefore not the collector but the broader managed-runtime
measurement path---bounds-checked array access, scheduler quanta, and coarser or
higher-overhead timing than the C collector. In short: Go's disadvantage is
real, reproducible, and a property of \emph{measurement}, not of the algorithm
or of the garbage collector specifically.\footnote{Absolute \texttt{min\_time}
values are far larger in Go than in C; whether that reflects timer units or true
overhead requires inspecting each collector's timing primitive, so we base the
cross-language claim on recovery reliability rather than on raw timing
magnitudes.}

\subsection{Dials: Signal Strength over Averaging, and a Design Lesson}
\label{sec:analysis-dial}
Within C at the 100k knee, additional \texttt{-repeat} does not help and the
largest \texttt{-evict-kb} helps most. This is consistent with the
averaging/signal split above: at a count near the threshold, the binding
constraint is the strength of the per-trace collision signal (improved by
stronger cache eviction), not additional per-trace averaging---and very large
\texttt{-repeat} may even lengthen each measurement enough to admit more
environmental perturbation. The Go dial sweep carries a methodological lesson
rather than a result: fixing the dial-sweep count at C's knee (100k) placed it
below Go's operating threshold ($\approx 265$k), so the sweep could not probe
Go's dial sensitivity. A per-implementation dial count, chosen relative to each
implementation's own threshold, would be required to compare dial effects across
languages.

\subsection{Positioning Against Prior Work}
\label{sec:analysis-priorwork}
The implemented attack follows the final-round cache-collision model of Bonneau
and Mironov, whose expanded final-round attack recovers a key from a median of
$2^{13}$--$2^{14.3}$ samples under cache eviction on 2006-era processors. Our C
threshold ($N_{50}\approx 2^{15.6}$; high-reliability $\approx 2^{16.2}$) is
roughly one order of magnitude higher. This is expected rather than
contradictory: the comparison mixes algorithm with platform and instrumentation.
Their figures use hardware cycle counters on a Pentium III/UltraSPARC, whereas
this laboratory records aggregate wall-clock time on an Apple M4 with a very
different cache hierarchy and timing resolution. The agreement to within an order
of magnitude, on entirely different hardware and with a coarser timer, indicates
that the leakage mechanism reproduces as predicted; the absolute count is a
property of this platform and instrument, not a universal constant.

\subsection{Threats to Validity}
\label{sec:analysis-threats}
\begin{itemize}
  \item \textbf{Single key.} All runs reuse one fixed AES key to isolate trace
        count; thresholds may vary across keys.
  \item \textbf{Sampling uncertainty.} Ten trials per count yield wide Wilson
        intervals, especially near 0/10 and 10/10; individual dial points should
        not be over-interpreted, and $N_{50}$ is an estimate, not a sharp value.
  \item \textbf{Run-to-run environment sensitivity.} Two clean C sweeps place the
        threshold near 50k and 100k respectively, and a contaminated sweep near
        150k; the threshold is sensitive to background load, which we control by
        the environment gate but cannot eliminate on a general-purpose machine.
  \item \textbf{Native-collection confound.} The cross-language comparison uses
        each language's own traces (different seeds), so it measures
        runtime-in-collection, not shared-dataset algorithmic consistency;
        the latter (draft RQ2) and stage-level processing time (RQ3) remain
        future work.
  \item \textbf{Reduced Go(GC-off) grid and platform confound.} The ablation uses
        a coarse count grid, and the prior-work comparison spans different
        hardware eras; both are stated where the numbers are used.
\end{itemize}
